% main.tex — Proyecto final: matriz SABSA para pagos inmediatos Bre-B.
% Compilar con XeLaTeX:  latexmk main.tex   (ver latexmkrc)
\documentclass[aspectratio=169,10pt]{beamer}
\usetheme[nosectionpages,accent=DA0921]{terminal}  % rojo oficial UR
\usepackage[spanish,es-noshorthands]{babel}

% caso.tex — datos del caso que cambian si se cambia de empresa.
% El contenido completo de la matriz vive en ../matriz/*.md (fuente del Excel);
% las slides lo resumen a mano.
\newcommand{\empresa}{NovaPay S.A.}
\newcommand{\empresaslug}{novapay}  % en los prompts de shell
\newcommand{\empresatipo}{compañía de financiamiento digital vigilada por la SFC}
\newcommand{\area}{Gerencia de Pagos Inmediatos}
\newcommand{\proyecto}{Vinculación a Bre-B como participante}
\newcommand{\autores}{Sergio Andrés Caceres Díaz\texorpdfstring{\\}{, } Over Haider Castrillón Valencia\texorpdfstring{\\}{, } Michael Steven Pinilla Mateus}
\newcommand{\autorescorto}{Caceres, Castrillón, Pinilla}  % pie de página


\title{Arquitectura de seguridad SABSA para pagos inmediatos Bre-B}
\subtitle{\empresa{} // \area}
\author[\autorescorto]{\autores}
\institute{Seguridad para Arquitectura Empresarial -- U. del Rosario}
\date{Noviembre 2026}

\terminalwindow{sabsa-breb.tex}
\terminalteamlabel{equipo}
\terminallogo{\includegraphics[height=2.6ex]{pics/ur-logo-white}}
\terminaltitlelogo{\includegraphics[height=1.25cm]{pics/ur-logo}}
\terminalfootprefix{sabsa}
\terminalsectionprompt{cd capa\_}
\terminalshell{\texttt{\$ ./sabsa -{}-empresa \empresaslug{} -{}-area bre-b\\
  > capas contextual..component ...... ok\\
  > matriz SABSA (xlsx) .............. ok}}

\begin{document}

\begin{frame}[plain,noframenumbering]
  \titlepage
\end{frame}

% ------------------------------------------------------------------ 01
\section{Misión}
\begin{frame}{Misión de la empresa}
  \prompt{cat /etc/\empresaslug/mision}
  \vskip8pt
  \begin{tbox}{> mision.txt}
    \large Facilitar que personas y comercios en Colombia muevan su dinero de
    forma \acc{inmediata}, \acc{segura} y \acc{sin costo} desde el celular,
    protegiendo la \acc{confianza} de los clientes como el activo principal de
    la compañía.
  \end{tbox}
  \vskip12pt
  \begin{columns}[T]
    \begin{column}{0.31\textwidth}
      \begin{pbox}\small\textbf{Quiénes:} \empresa, \empresatipo.\end{pbox}
    \end{column}
    \begin{column}{0.31\textwidth}
      \begin{pbox}\small\textbf{Qué:} Cuentas digitales, transferencias y pagos QR por Bre-B.\end{pbox}
    \end{column}
    \begin{column}{0.31\textwidth}
      \begin{pbox}\small\textbf{Por qué:} Bre-B es interoperable, 24/7 y gratuito, una oportunidad para captar a quienes usan efectivo.\end{pbox}
    \end{column}
  \end{columns}
\end{frame}

% ------------------------------------------------------------------ 02
\section{Área}
\begin{frame}{Objetivo del área: Pagos Inmediatos Bre-B}
  \begin{tbox}{\texttt{> objetivo -{}-area bre-b}}
    Operar las transferencias y pagos inmediatos de \empresa{} sobre Bre-B
    (\acc{24/7}, liquidación en \acc{menos de 20 s}), administrando las llaves
    de los clientes ante el Directorio Centralizado (DICE) y la liquidación en
    el MOL del Banco de la República, con fraude y disponibilidad dentro de los
    umbrales de la Junta.
  \end{tbox}
  \vskip10pt
  \begin{columns}[T]
    \begin{column}{0.48\textwidth}
      \begin{tbox}{alcance}
        \footnotesize
        \begin{cellitems}
          \item Inscripción y portabilidad de llaves (celular, cédula, correo, alfanumérica).
          \item Envío, recepción y pagos QR a comercios.
          \item Conciliación con el MOL, disputas y fraude.
        \end{cellitems}
      \end{tbox}
    \end{column}
    \begin{column}{0.48\textwidth}
      \begin{tbox}{metas (KPI)}
        \footnotesize
        \begin{cellitems}
          \item Disponibilidad $\geq$ 99,95\,\% mensual.
          \item Transacción extremo a extremo $<$ 20\,s.
          \item Fraude $<$ 5 pb del valor transado.
          \item Cero incidentes de fuga de datos personales.
        \end{cellitems}
      \end{tbox}
    \end{column}
  \end{columns}
\end{frame}

% ------------------------------------------------------------------ 03
\section{Contextual}
\begin{frame}{Contextual \textmd{// requisitos de negocio}}
  \prompt{sabsa --layer contextual --view business --owner "Junta, CEO, VP Pagos, CISO"}
  \vskip4pt
  \begin{sabsagrid}
    \sabsacell{WHAT}{activos}{\begin{cellitems}
      \item Servicio 24/7 y $<$ 20\,s.
      \item Fondos y datos de clientes.
      \item Reputación y confianza.
    \end{cellitems}}
    \sabsacell{WHY}{amenazas}{\begin{cellitems}
      \item Fraude APP, ATO (SIM swap, phishing), mulas.
      \item DDoS y caída de nube o enlace.
      \item Insiders, supply chain, ransomware.
    \end{cellitems}}
    \sabsacell{HOW}{procesos}{\begin{cellitems}
      \item Onboarding / KYC.
      \item Gestión de llaves ante el DICE.
      \item Envío, recepción, QR, conciliación.
    \end{cellitems}}
    \sabsacell{WHO}{organización}{\begin{cellitems}
      \item Junta, CEO, VP Pagos, CISO.
      \item Fraude, Riesgos (SARO), Datos.
      \item BanRep, nube, biometría, telcos.
    \end{cellitems}}
    \sabsacell{WHERE}{sitios}{\begin{cellitems}
      \item Bogotá: operación y SOC.
      \item Nube: región principal + alterna.
      \item Enlace dedicado a BanRep.
    \end{cellitems}}
    \sabsacell{WHEN}{tiempo}{\begin{cellitems}
      \item Sin ventana de cierre.
      \item Picos 5$\times$ en quincena y primas.
      \item Reporte de incidentes a la SFC.
    \end{cellitems}}
  \end{sabsagrid}
  \src{SFC CE 007/2018 (hoy en la CBJ, CE 006/2025) · BanRep CRE DSP-465/2025}
\end{frame}

% ------------------------------------------------------------------ 04
\section{Conceptual}
\begin{frame}{Conceptual \textmd{// estrategias de seguridad}}
  \prompt{sabsa --layer conceptual --view architect --owner "CISO"}
  \vskip4pt
  \begin{sabsagrid}
    \sabsacell{WHAT}{atributos}{\begin{cellitems}
      \item Disponible, confidencial, íntegro.
      \item Autenticado, no repudiable.
      \item Resistente al fraude, cumplido.
    \end{cellitems}}
    \sabsacell{WHY}{control}{\begin{cellitems}
      \item Sin MFA + dispositivo, no hay pago.
      \item Todo pago pasa por scoring.
      \item Cifrado con llaves en HSM.
    \end{cellitems}}
    \sabsacell{HOW}{estrategias}{\begin{cellitems}
      \item \acc{Zero Trust}.
      \item Defensa en profundidad.
      \item Privacy by design, DevSecOps.
    \end{cellitems}}
    \sabsacell{WHO}{dueños}{\begin{cellitems}
      \item Proceso: VP Pagos.
      \item Riesgo ciber: CISO; fraude: Gte.
      \item Custodios: Plataforma y Datos.
    \end{cellitems}}
    \sabsacell{WHERE}{dominios}{\begin{cellitems}
      \item Cliente · Exposición (DMZ).
      \item Core de pagos · Datos sensibles.
      \item Interbancario · Gestión.
    \end{cellitems}}
    \sabsacell{WHEN}{riesgo}{\begin{cellitems}
      \item NIST CSF 2.0 + SARO.
      \item Revisión trimestral.
      \item KRI: fraude, MTTD, MTTR.
    \end{cellitems}}
  \end{sabsagrid}
  \src{ISO/IEC 27001:2022 · NIST CSF 2.0}
\end{frame}

% ------------------------------------------------------------------ 05
\section{Logical}
\begin{frame}{Logical \textmd{// servicios de seguridad}}
  \prompt{sabsa --layer logical --view designer --owner "Arquitecto de Seguridad"}
  \vskip4pt
  \begin{sabsagrid}
    \sabsacell{WHAT}{información}{\begin{cellitems}
      \item Cliente, KYC, biometría.
      \item Llave $\leftrightarrow$ cuenta (DICE).
      \item Transacciones, secretos, logs.
    \end{cellitems}}
    \sabsacell{WHY}{políticas}{\begin{cellitems}
      \item MFA fuerte y límites por riesgo.
      \item Bloqueo 24\,h tras cambio de SIM.
      \item TLS 1.3, mTLS y AES-256.
    \end{cellitems}}
    \sabsacell{HOW}{servicios}{\begin{cellitems}
      \item CIAM y autorización central.
      \item Scoring de fraude en línea.
      \item Cripto, monitoreo y auditoría.
    \end{cellitems}}
    \sabsacell{WHO}{confianza}{\begin{cellitems}
      \item Cliente: solo sus llaves.
      \item Agente: datos enmascarados.
      \item Ingeniero: just-in-time.
    \end{cellitems}}
    \sabsacell{WHERE}{flujos}{\begin{cellitems}
      \item Cliente $\to$ API GW + WAF.
      \item Core $\to$ datos: solo servicio.
      \item Core $\to$ BanRep: mTLS dedicado.
    \end{cellitems}}
    \sabsacell{WHEN}{vigencias}{\begin{cellitems}
      \item Sesión 5\,min; token 10\,min.
      \item Scoring $<$ 300\,ms.
      \item Certificados cada 90 días.
    \end{cellitems}}
  \end{sabsagrid}
\end{frame}

% ------------------------------------------------------------------ 06
\section{Physical}
\begin{frame}{Physical \textmd{// mecanismos de seguridad}}
  \prompt{sabsa --layer physical --view builder --owner "Dir. Plataforma Cloud"}
  \vskip4pt
  \begin{sabsagrid}
    \sabsacell{WHAT}{datos}{\begin{cellitems}
      \item Restringido: tokenizado.
      \item Confidencial: cifrado.
      \item Catálogo con dueño y retención.
    \end{cellitems}}
    \sabsacell{WHY}{reglas}{\begin{cellitems}
      \item Reglas de fraude (velocidad, mula).
      \item Playbooks ATO, DDoS, ransomware.
      \item Parches críticos en 7 días.
    \end{cellitems}}
    \sabsacell{HOW}{mecanismos}{\begin{cellitems}
      \item Passkeys FIDO2 + device binding.
      \item WAF, anti-DDoS, service mesh.
      \item HSM/KMS, tokenización, SIEM.
    \end{cellitems}}
    \sabsacell{WHO}{acceso}{\begin{cellitems}
      \item SSO + llave FIDO2 interna.
      \item RBAC + ABAC (monto, horario).
      \item PAM con sesiones grabadas.
    \end{cellitems}}
    \sabsacell{WHERE}{red}{\begin{cellitems}
      \item VPC por ambiente.
      \item Subred por dominio.
      \item 3 zonas + región alterna.
    \end{cellitems}}
    \sabsacell{WHEN}{tiempos}{\begin{cellitems}
      \item \acc{RTO 15\,min, RPO 0}.
      \item Failover de zona $<$ 1\,min.
      \item Backups inmutables diarios.
    \end{cellitems}}
  \end{sabsagrid}
\end{frame}

% ------------------------------------------------------------------ 07
\section{Component}
\begin{frame}{Component \textmd{// productos y tecnologías}}
  \prompt{sabsa --layer component --view tradesman --owner "Líder SecOps"}
  \vskip4pt
  \begin{sabsagrid}
    \sabsacell{WHAT}{plataforma}{\begin{cellitems}
      \item Amazon EKS, Aurora, MSK.
      \item AWS CloudHSM y KMS.
      \item HashiCorp Vault.
    \end{cellitems}}
    \sabsacell{WHY}{riesgo}{\begin{cellitems}
      \item Feedzai (fraude, ML).
      \item Microsoft Sentinel (SIEM/SOAR).
      \item Tenable, Snyk, Security Hub.
    \end{cellitems}}
    \sabsacell{HOW}{protocolos}{\begin{cellitems}
      \item ISO 20022, OIDC, WebAuthn.
      \item Cloudflare, Istio (mTLS).
      \item CI: Semgrep, Trivy, ZAP.
    \end{cellitems}}
    \sabsacell{WHO}{identidades}{\begin{cellitems}
      \item Clientes: Auth0 / Keycloak.
      \item Empleados: Entra ID + YubiKey.
      \item PAM: CyberArk / Teleport.
    \end{cellitems}}
    \sabsacell{WHERE}{nodos}{\begin{cellitems}
      \item EKS prod (3 AZ) + DR.
      \item Sin IP pública en core/datos.
      \item Direct Connect a BanRep.
    \end{cellitems}}
    \sabsacell{WHEN}{agenda}{\begin{cellitems}
      \item Pentest trimestral.
      \item Escaneo semanal y por build.
      \item DRP semestral; NTP en todo.
    \end{cellitems}}
  \end{sabsagrid}
\end{frame}

% ------------------------------------------------------------------ 08
\section{Conclusiones}
\begin{frame}{Conclusiones}
  \begin{cellitems}\small
    \item \textbf{Trazabilidad:} Cada producto de la capa Component se justifica
      hacia arriba hasta un requisito de negocio de Bre-B, de forma que ningún
      control queda sin dueño ni sin motivo.
    \vskip4pt
    \item \textbf{Fraude antes que intrusión:} En pagos inmediatos el dinero no
      se puede reversar, por lo cual la autenticación fuerte, el device binding
      y el scoring en línea pesan tanto como el perímetro.
    \vskip4pt
    \item \textbf{Disponibilidad como requisito de seguridad:} Un servicio 24/7
      sin ventana de cierre obliga a diseñar la resiliencia (RTO 15\,min,
      RPO 0) desde la capa conceptual y no como un añadido.
    \vskip4pt
    \item \textbf{Cumplimiento integrado:} La SFC, el Banco de la República y la
      Ley 1581 se tratan como atributos de negocio y no como un checklist al
      final del proyecto.
  \end{cellitems}
  \vskip8pt
  \begin{ibox}SABSA = seguridad derivada del negocio, no agregada a la tecnología\end{ibox}
\end{frame}

\terminalscreen{\texttt{\$ sabsa -{}-preguntas}}{¿PREGUNTAS?}{\empresa{} \quad//\quad \area}

\end{document}
